\textbf{Phần 1. Trắc nghiệm lựa chọn}
\setcounterpageref{ex}{0}

\begin{ex} %[3H3V51]
Xăng-ti-mét vuông là diện tích của hình vuông có cạnh dài
\choice
{\True $1\text{ cm}$}
{$1\text{ dm}$}
{$1\text{ m}$}
{$1\text{ mm}$}
\loigiai{
Xăng-ti-mét vuông là diện tích của hình vuông có cạnh dài $1\text{ cm}$, kí hiệu là $\text{cm}^2$.
\textbf{Đáp án của câu hỏi là A}
}
\end{ex}

\begin{ex} %[3H3V51]
Một hình được ghép bởi 15 ô vuông nhỏ, mỗi ô vuông có diện tích $1\text{ cm}^2$. Diện tích của hình đó là
\choice
{$30\text{ cm}^2$}
{$60\text{ cm}^2$}
{\True $15\text{ cm}^2$}
{$20\text{ cm}^2$}
\loigiai{
Diện tích cả hình: $1 \times 15 = 15\text{ cm}^2$.
\textbf{Đáp án của câu hỏi là C}
}
\end{ex}

\begin{ex} %[3H3V51]
Hình $A$ có diện tích $28\text{ cm}^2$, hình $B$ có diện tích $32\text{ cm}^2$. Khẳng định nào sau đây là đúng?
\choice
{Diện tích hình $A$ lớn hơn hình $B$}
{\True Diện tích hình $B$ lớn hơn diện tích hình $A$ là $4\text{ cm}^2$}
{Diện tích hai hình bằng nhau}
{Không thể so sánh được diện tích hai hình}
\loigiai{
$32 - 28 = 4\text{ cm}^2$. Hình $B$ lớn hơn hình $A$ là $4\text{ cm}^2$.
\textbf{Đáp án của câu hỏi là B}
}
\end{ex}

\begin{ex} %[3H3V51]
Người ta cắt một tấm bìa có diện tích $50\text{ cm}^2$ thành 2 mảnh: mảnh thứ nhất có diện tích $28\text{ cm}^2$. Mảnh bìa thứ hai có diện tích là
\choice
{$24\text{ cm}^2$}
{$20\text{ cm}^2$}
{$18\text{ cm}^2$}
{\True $22\text{ cm}^2$}
\loigiai{
Diện tích mảnh thứ hai: $50 - 28 = 22\text{ (cm}^2)$.
\textbf{Đáp án của câu hỏi là D}
}
\end{ex}

\begin{ex} %[3H3V51]
Một hình phẳng gồm 4 hàng ô vuông, mỗi hàng có 6 ô vuông nhỏ cạnh $1\text{ cm}$. Diện tích của hình phẳng đó là
\choice
{\True $24\text{ cm}^2$}
{$20\text{ cm}^2$}
{$28\text{ cm}^2$}
{$36\text{ cm}^2$}
\loigiai{
Tổng số ô vuông nhỏ: $6 \times 4 = 24$ ô. Diện tích là $24\text{ cm}^2$.
\textbf{Đáp án của câu hỏi là A}
}
\end{ex}

\begin{ex} %[3H3V51]
Tính: $35\text{ cm}^2 + 45\text{ cm}^2 : 5$. Kết quả là
\choice
{$16\text{ cm}^2$}
{$40\text{ cm}^2$}
{\True $44\text{ cm}^2$}
{$50\text{ cm}^2$}
\loigiai{
$45 : 5 = 9$; $35 + 9 = 44\text{ cm}^2$.
\textbf{Đáp án của câu hỏi là C}
}
\end{ex}

\begin{ex} %[3H3V51]
Một tờ giấy màu đỏ có diện tích $80\text{ cm}^2$. Tờ giấy màu xanh có diện tích bằng $\dfrac{1}{4}$ tờ giấy màu đỏ. Tổng diện tích cả hai tờ giấy màu là
\choice
{$120\text{ cm}^2$}
{\True $100\text{ cm}^2$}
{$90\text{ cm}^2$}
{$110\text{ cm}^2$}
\loigiai{
Giấy xanh: $80 : 4 = 20\text{ cm}^2$. Tổng: $80 + 20 = 100\text{ cm}^2$.
\textbf{Đáp án của câu hỏi là B}
}
\end{ex}

\begin{ex} %[3H3V51]
Có 5 miếng bìa hình chữ nhật giống nhau, mỗi miếng bìa có diện tích $18\text{ cm}^2$. Ghép 5 miếng bìa đó lại không đè lên nhau ta được hình mới có diện tích là
\choice
{$80\text{ cm}^2$}
{$85\text{ cm}^2$}
{$95\text{ cm}^2$}
{\True $90\text{ cm}^2$}
\loigiai{
Diện tích hình mới: $18 \times 5 = 90\text{ (cm}^2)$.
\textbf{Đáp án của câu hỏi là D}
}
\end{ex}

\begin{ex} %[3H3V51]
Diện tích của một con tem thư thường vào khoảng
\choice
{\True $6\text{ cm}^2$}
{$60\text{ cm}^2$}
{$6\text{ m}^2$}
{$600\text{ cm}^2$}
\loigiai{
Con tem thư nhỏ kích thước khoảng $2\text{ cm} \times 3\text{ cm} = 6\text{ cm}^2$.
\textbf{Đáp án của câu hỏi là A}
}
\end{ex}

\begin{ex} %[3H3V51]
Tính giá trị biểu thức: $120\text{ cm}^2 - 15\text{ cm}^2 \times 4$.
\choice
{$50\text{ cm}^2$}
{$70\text{ cm}^2$}
{\True $60\text{ cm}^2$}
{$80\text{ cm}^2$}
\loigiai{
$15 \times 4 = 60$; $120 - 60 = 60\text{ cm}^2$.
\textbf{Đáp án của câu hỏi là C}
}
\end{ex}

\begin{ex} %[3H3V51]
Một hình chữ nhật được chia làm 2 hình vuông bằng nhau, mỗi hình vuông có diện tích $25\text{ cm}^2$. Diện tích hình chữ nhật đó là
\choice
{$40\text{ cm}^2$}
{\True $50\text{ cm}^2$}
{$100\text{ cm}^2$}
{$75\text{ cm}^2$}
\loigiai{
Diện tích hình chữ nhật: $25 \times 2 = 50\text{ cm}^2$.
\textbf{Đáp án của câu hỏi là B}
}
\end{ex}

\begin{ex} %[3H3V51]
Hình nào có diện tích lớn nhất trong các hình sau?
\choice
{Hình có diện tích $20\text{ cm}^2$}
{Hình gồm 20 ô vuông cạnh $1\text{ cm}$}
{Hình có diện tích $24\text{ cm}^2$}
{\True Hình gồm 25 ô vuông cạnh $1\text{ cm}$}
\loigiai{
Hình gồm 25 ô vuông có diện tích $25\text{ cm}^2$ là lớn nhất.
\textbf{Đáp án của câu hỏi là D}
}
\end{ex}

\textbf{Phần 2. Trắc nghiệm đúng sai}
\setcounterpageref{ex}{0}

\begin{ex} %[3H3V51]
Xét tính đúng sai của các nhận định về xăng-ti-mét vuông:
\choiceTF
{\True Xăng-ti-mét vuông kí hiệu là $\text{cm}^2$.}
{$1\text{ cm}^2$ là độ dài một đoạn thẳng dài $1\text{ cm}$.}
{\True Khi ghép các hình phẳng lại không chồng lên nhau thì diện tích hình mới bằng tổng diện tích các hình thành phần.}
{Hai hình có diện tích bằng nhau thì luôn có hình dạng giống hệt nhau.}
\loigiai{
\begin{itemize}
    \item Ý a) Đúng: Kí hiệu chuẩn.
    \item Ý b) Sai: Sai vì xăng-ti-mét vuông là đơn vị đo diện tích, không phải độ dài.
    \item Ý c) Đúng: Tính chất cộng diện tích.
    \item Ý d) Sai: Sai vì một hình chữ nhật và một hình vuông có thể cùng diện tích.
\end{itemize}
\begin{itemize}
    \item a) \textbf{Đúng}.
    \item b) \textbf{Sai}.
    \item c) \textbf{Đúng}.
    \item d) \textbf{Sai}.
\end{itemize}
}
\end{ex}

\begin{ex} %[3H3V51]
Một miếng bìa lớn có diện tích $96\text{ cm}^2$. Người ta cắt miếng bìa thành các mảnh nhỏ diện tích $8\text{ cm}^2$:
\choiceTF
{\True Cắt được tất cả 12 mảnh bìa nhỏ.}
{\True Tổng diện tích của 5 mảnh bìa nhỏ là $40\text{ cm}^2$.}
{Sau khi cắt ra 12 mảnh nhỏ thì diện tích bị hao hụt một nửa.}
{\True Mỗi mảnh bìa nhỏ có diện tích gấp 8 lần ô vuông $1\text{ cm}^2$.}
\loigiai{
\begin{itemize}
    \item Ý a) Đúng: $96 : 8 = 12$ mảnh.
    \item Ý b) Đúng: $8 \times 5 = 40\text{ cm}^2$.
    \item Ý c) Sai: Sai vì tổng diện tích vẫn bảo toàn.
    \item Ý d) Đúng: Đúng theo định nghĩa.
\end{itemize}
\begin{itemize}
    \item a) \textbf{Đúng}.
    \item b) \textbf{Đúng}.
    \item c) \textbf{Sai}.
    \item d) \textbf{Đúng}.
\end{itemize}
}
\end{ex}

\textbf{Phần 3. Trả lời ngắn}
\setcounterpageref{ex}{0}

\begin{ex} %[3H3V51]
Một hình được ghép bởi 18 ô vuông nhỏ cạnh $1\text{ cm}$. Tính diện tích hình đó theo xăng-ti-mét vuông.
\shortans{18}
\loigiai{
\textbf{Đáp số:} $18$\\
Diện tích là: $18\text{ cm}^2$.
}
\end{ex}

\begin{ex} %[3H3V51]
Tính giá trị: $48\text{ cm}^2 + 36\text{ cm}^2$ theo xăng-ti-mét vuông.
\shortans{84}
\loigiai{
\textbf{Đáp số:} $84$\\
$48 + 36 = 84\text{ (cm}^2)$.
}
\end{ex}

\begin{ex} %[3H3V51]
Một tấm bìa diện tích $100\text{ cm}^2$ bị cắt đi một góc có diện tích $15\text{ cm}^2$. Phần bìa còn lại có diện tích bao nhiêu xăng-ti-mét vuông?
\shortans{85}
\loigiai{
\textbf{Đáp số:} $85$\\
$100 - 15 = 85\text{ (cm}^2)$.
}
\end{ex}

\begin{ex} %[3H3V51]
Có 4 hình vuông nhỏ, mỗi hình có diện tích $9\text{ cm}^2$. Ghép lại được hình lớn có diện tích bao nhiêu xăng-ti-mét vuông?
\shortans{36}
\loigiai{
\textbf{Đáp số:} $36$\\
$9 \times 4 = 36\text{ (cm}^2)$.
}
\end{ex}

\begin{ex} %[3H3V51]
Một hình chữ nhật có diện tích $72\text{ cm}^2$ được chia thành 8 phần bằng nhau. Mỗi phần có diện tích bao nhiêu xăng-ti-mét vuông?
\shortans{9}
\loigiai{
\textbf{Đáp số:} $9$\\
$72 : 8 = 9\text{ (cm}^2)$.
}
\end{ex}

\begin{ex} %[3H3V51]
Tính: $15\text{ cm}^2 \times 6 - 25\text{ cm}^2$.
\shortans{65}
\loigiai{
\textbf{Đáp số:} $65$\\
$90 - 25 = 65\text{ (cm}^2)$.
}
\end{ex}
